\chapterrule
\section{用 MLP 的眼光重新理解记忆检索}
\phantomsection\label{intro:09-nonlinearity-and-memory}
\chaptercompass{匹配分数为什么需要非线性？}{先把 Key/Value 堆成矩阵，再用非线性筛选匹配分数，按权重组合每条 Value。}{$xK^\top:1\times T_k\to\sigma(\cdot)V:1\times d_v$}{在 MLP 解释中 K/V 对应可训练的层权重}
\modelpipelinefocus{3}{非线性记忆检索}

上一章的两种线性写法都只在写入端减少串扰，读取时所有非零匹配仍按原值参与输出。
本章转到读取端：先把多条 Key、Value 堆成矩阵，再让权重随匹配强弱非线性变化。

\subsection{把 Key 和 Value 堆成矩阵：从匹配到读出}
\phantomsection\label{intro:08-stacking-key-value}
设 Query $x:1\times d_k$，$T_k$ 条 Key 按行排成
$K:T_k\times d_k$，Value 按行排成 $V:T_k\times d_v$。
一次查询经过三步：
\begin{equation}
\boxed{\underbrace{y}_{1\times d_v}
=\underbrace{\sigma(xK^\top)}_{1\times T_k}
\underbrace{V}_{T_k\times d_v}.}
\end{equation}
这与前面 $\softmax(QK^\top)V$ 采用同一个行向量约定。

\subsubsection{两条二维记忆，手算一次}
\localfixture{二维 Key/Value；独立于三维主例}
\[
x=(1,-1),\quad
K=\begin{bmatrix}1&0\\0.6&0.8\end{bmatrix},\quad
V=\begin{bmatrix}2&1\\0&3\end{bmatrix}.
\]
两次点积分别是 $1$ 和 $-0.2$，
合并成一行：
\[
xK^\top=(1,-1)\begin{bmatrix}1&0.6\\0&0.8\end{bmatrix}
=\boxed{(1,-0.2)}.
\]
先逐元素 ReLU，再按行混合 Value：
\begin{example}
\[
a=\relu(xK^\top)=(1,0),\qquad
aV=(1,0)\begin{bmatrix}2&1\\0&3\end{bmatrix}
=\boxed{(2,1)}.
\]
第二条记忆的负匹配被筛掉，因此未参与输出。
\end{example}
ReLU 是作用于分数，而不是事先作用于 Key 矩阵。

\subsubsection{从二维推广到不同维度}
$T_k$ 是记忆条数，$d_k$ 是每条 Key 的宽度，$d_v$ 是 Value 的宽度。
\[
(1\times d_k)(d_k\times T_k)(T_k\times d_v)=(1\times d_v).
\]
例如 $T_k=100$、$d_k=64$、$d_v=32$，
匹配结果是 $1\times100$，输出是 $1\times32$。
每个 Token 占一行的规则始终不变。

\subsubsection{插入非线性究竟改变了什么？}
如果不插入 ReLU，则
\[
\mathcal{M}=K^\top V=\sum_i k_i^\top v_i,
\qquad x\mathcal{M}=(xK^\top)V.
\]
这次例子的线性读出是
\[
x\mathcal{M}=(1,-0.2)V=\boxed{(2,0.4)},
\]
而非线性读出为 $\relu(xK^\top)V=(2,1)$。
\textbf{不插 ReLU 时，两步可预先合并成固定状态 $\mathcal{M}=K^\top V$；
插入 ReLU 后，每条匹配先按正负筛选，读出只能逐条匹配后再组合。}
Softmax Attention 也先逐条匹配、归一化再读出，属于后一类。

\noindent{\sffamily\footnotesize\color{blue}\textbf{［复算］}}
运行 \texttt{python3 scripts/memory\_demo.py}，
检查两种读出结果、存储形状与串扰。

\begin{idea}
统一记住一条式子：$\sigma(xK^\top)V$。
先匹配所有 Key，再把这些分数当作读取各行 Value 的权重。
\end{idea}
